% !TeX root = ../main.tex
\section{Introduction}\label{sec:intro}

Modern NVMe SSDs deliver high IOPS and throughput only when requests are issued concurrently. Like
memory, storage shows a queueing-driven throughput--latency tradeoff: deeper queues raise device
utilization up to a saturation knee, and past it latency climbs sharply for little extra throughput.
This project characterizes that behaviour on a Micron 2210 \qty{512}{\giga\byte} QLC NVMe SSD. It
measures QD=1 latency, block-size and read/write-mix effects, the queue-depth knee, tail latency and
device state (working set, SLC-cache exhaustion, temperature). All I/O is \code{O_DIRECT} against
disposable files. The primary engine is a small \code{libaio} benchmark, and FIO cross-checks it
(\cref{sec:method,sec:validation}).

This report answers:
\begin{enumerate}
  \item What are the QD=1 latencies for small random and large sequential I/O? (\cref{sec:zero-queue})
  \item How do block size, access pattern, read/write mix and queue depth shape IOPS, MiB/s and latency,
        and where is the PCIe interface ceiling? (\cref{sec:block-size,sec:rw-mix})
  \item Where is the queue-depth knee, by what criterion, and what does Little's Law say about the
        harness? (\cref{sec:queue-depth})
  \item What do tail latency and device state reveal: working-set size, SLC-cache exhaustion and steady
        state, data layout, and temperature? (\cref{sec:tail,sec:device-state})
  \item Which setup mistake produces misleading numbers? (\cref{sec:pitfall})
\end{enumerate}

In summary, QD=1 random reads cost
\csvq[1]{zero_queue}{mode=read,pattern=rand}{lat_mean_us_median}\,\si{\micro\second} and writes only
\csvq[1]{zero_queue}{mode=write,pattern=rand}{lat_mean_us_median}\,\si{\micro\second}, because the write
buffer absorbs them. With one request in flight, sequential reads are limited by per-request latency
rather than the PCIe~3.0~$\times$4 link. Deeper queues raise them to only \qty{57}{\percent} of the link,
where the NAND back end becomes the limit (\cref{sec:block-size}). The random-read knee sits at QD~\csvq{queue_depth_knee}{mode=read}{knee_qd}.
Tails stretch with queueing and with background flash work. Once the SLC cache is full, sustained writes drop
$\ssdburstratio\times$ to the direct-to-QLC rate, and the drive's recent write history changes both read tails and
sequential-read speed (\cref{sec:device-state}). Buffered I/O over-reports throughput
by more than an order of magnitude because it measures the page cache, not the device.
